(1) 介质间的球形高斯面给 $4\pi r^2D_r=Q$，故 $a<r<b$ 时 $\mathbf D=Q\hat r/(4\pi r^2)$、$\mathbf E=Q\hat r/(4\pi\varepsilon r^2)$；理想导体金属内部静电场为零，且因总电荷为零、球对称，外部 $r>b$ 的电场也为零。
(2) 两导体势差 $\Delta V=\Phi(a)-\Phi(b)=\int_a^b E_r\,dr=Q(1/a-1/b)/(4\pi\varepsilon)$。由 $C=Q/\Delta V$ 得 $C=4\pi\varepsilon ab/(b-a)$，量纲为法拉，且 $b\to a^+$ 时电容增大，符合窄间隙趋势。
(3) 场能积分 $U=\frac12\int_a^b\varepsilon E^2\,4\pi r^2dr=Q^2(1/a-1/b)/(8\pi\varepsilon)=Q^2/(2C)$。在 $Q$ 固定时增大 $\varepsilon$ 使电场强度、场能按 $\varepsilon^{-1}$ 降低，因此介质极化会减小固定自由电荷系统的储能。
(4) 恒压时总电荷 $Q=CV=4\pi\varepsilon abV/(b-a)$，场能 $U_{\rm field}=\frac12CV^2$。电源对充电过程做功 $\int V\,dQ=VQ$（恒压），故用于求静态介质或导体机械力的有效能为 $U_{\rm field}-VQ=-CV^2/2$；若只求场能梯度而忽略电源，力的符号可能错误。
